% Create a "code name" for each figure. In the maintext, when you want that figure to appear in the manuscript call that function (e.g. \figCodeName).

%%%% Create and format Figure with subpanels: %%%%
% \generateFigSubpanels[optionalRotationDegrees]{figCodeName}{figureNamePath}{size:propOfTextwidth}
%     {main caption}
%     {{
%         {subcaption1},
%         {subcaption2},
%         {subcaption3}
%     }}

% Then in the reference your figures with \cref{fig:figCodeName} and Reference subpanels with \cref{fig:figCodeName:A}


%%%% Function to create and format figures without subpanels %%%%

% \generateFig[optionalRotationDegrees]{figCodeName}{figureNamePath}{size:propOfTextwidth}
%     {Bold text main caption}
%     {Smaller normal text caption}

% Then in the main text reference your figures with \cref{fig:figCodeName}

% You can also use \generateSidwaysFigSubpanels or \generateSidewaysFig to rotate the caption with the figure.

%%%% Example figure with codename: figureSuggestions %%%%
\generateFig{schematic}
{main_figs/Figure1.pdf}{1}
    {Schematic of the workflow of the study.}
    {We analyzed genome-wide methylation and genotyping data from blood samples from 621 AD and non-demented control individuals from the MAGENTA study. To replicate our findings, we analyzed data from 1,394 healthy African Americans from the GENOA study, 422 healthy African Americans from the Grady Trauma Project, and 729 healthy Whites from the NSPHS study. We applied a set of first-, second-, and third-generation methylation clocks to the individuals and estimated their genetic ancestry. This enabled us to explore the performance of methylation clocks in individuals with different genetic ancestries.}
    
\generateFig{correlationsancestry}
{main_figs/Figure2.pdf}{1}
    {Methylation clock accuracy is lower in cohorts with substantial African genetic ancestry.}
    {{
        {\textbf{A}: Pearson correlation between chronological age and DNAm age predicted by the Horvath clock for controls in the White MAGENTA cohort. The correlation of $r = 0.72$ is consistent with previous studies of similar age groups.}
        {\textbf{B}: Pearson correlation between chronological age and DNAm age predicted by the Horvath clock for the genetically admixed cohorts in MAGENTA. For each cohort, the adjacent boxplots display the distribution of global ancestry proportions: European (CEU, red), African (YRI, green), and Amerindigenous (PEL, blue). The cohorts with substantial African ancestry---African Americans and Puerto Ricans---exhibit lower correlations ($r = 0.51$ and $0.45$, respectively) compared to the Cubans ($r = 0.68$) and Peruvians ($r = 0.72$).}
        {\textbf{C}: Relative accuracy of the Horvath clock across cohorts compared to Whites. The bar plot shows the difference in Pearson correlation coefficients relative to the White cohort baseline ($r = 0.72$). Asterisks indicate a statistically significant difference from the baseline (* \textit{p} $<$ 0.05).}
    }}

\generateFig{magentareplication}
{main_figs/Figure3.pdf}{1}
    {Lower methylation clock accuracy in African Americans holds across cohorts.}
    {{
        Scatter plots of chronological age versus Horvath DNAm age in three independent replication cohorts: Swedish Whites (NSPHS), GENOA Study African Americans, and Grady Trauma Project African Americans.
        \textbf{Top Row}: Age predictions for the full age range of each cohort. The Swedish White cohort exhibits the highest accuracy ($r = 0.97$, MAE = 3.04), while both African American cohorts show lower correlations ($r = 0.85$ and $r = 0.88$) and higher error rates (MAE = 3.83 and 5.14, respectively).
        \textbf{Bottom Row}: Age predictions restricted to individuals $\ge$ 55 years old, similar to the demographics of the MAGENTA cohorts. In this age-restricted subset, the disparity in clock performance is consistent, with the Swedish White cohort maintaining a significantly higher correlation ($r = 0.80$) compared to the GENOA ($r = 0.67$) and Grady ($r = 0.49$) African American cohorts.
    }}

\generateFig{accuracymultipleclocks}
{main_figs/Figure4.pdf}{1}
    {Lower methylation clock accuracy in cohorts with African ancestry holds across multiple clocks.}
    {{
        Pearson correlation coefficients between chronological age and DNAm age predicted by the \textbf{Horvath} (top), \textbf{Hannum} (middle), and \textbf{Zhang\_EN} (bottom) clocks across all MAGENTA and replication (GENOA, Grady, Swedish) cohorts.
        The vertical dashed line in each panel represents the baseline correlation observed in the MAGENTA White cohort.
        Cohorts with substantial African ancestry (African American and Puerto Rican, red bars) consistently exhibit lower correlations compared to Whites, Peruvians, and Cubans across all three clock models.
        Asterisks indicate a statistically significant difference in correlation compared to the White MAGENTA cohort baseline (* \textit{p} $<$ 0.05).
    }}

\generateFig{alzheimersanalysis}
{main_figs/Figure5.pdf}{1}
    {Methylation clocks do not consistently identify accelerated aging in admixed Alzheimer's cohorts.}
    {{
        {\textbf{A}: Comparison of the distributions of Horvath intrinsic age acceleration for AD patients and non-demented controls for each of the admixed cohorts in MAGENTA. AD patients do not show significantly higher age acceleration in any of the admixed cohorts. In contrast, the AD cases had significantly greater acceleration than controls in the white cohort (\textbf{Supplementary Figure 6}). NS = Not significant.}
        {\textbf{B}: Median differences (in years) in intrinsic age acceleration between AD patients and non-demented controls for five methylation clocks for each cohort in MAGENTA. The clocks do not consistently identfy accelerated aging in AD across cohorts, and the results also vary within cohorts. * \textit{p} $<$ 0.05.}
    }}

% \generateFig{meqtlsgnomAD}
% {%main_figs/Figure5.pdf}{0.53}
%     {Common meQTLs affect most Horvath clock CpGs and vary in frequency across ancestries.}
%     {
%         {\textbf{A}: The allele frequency distribution of the 29,033 unique variants associated with methylation levels at Horvath clock CpGs. Allele frequencies were computed over the 76,215 genomes in gnomAD version 4.1. Inset: Out of the 353 CpGs in the Horvath clock, 271 (77\%) have at least one meQTL, i.e., a genetic variant that is associated with methylation level.}
%         {\textbf{B}: Clock meQTL have significantly higher alelle frequency in individuals with African genetic ancestry from gnomAD than all other ancestry groups (median 0.068 for African vs. 0.004--0.046; p $<$ $3.85 \times 10^{-25}$). }
%         {\textbf{C}: Clock meQTL have significantly higher allele frequency on African local ancestry genomic segments in 7,612 Latino admixed individuals with varying proportions of European, American, and African ancestry from gnomAD. Inset: The distribution of difference in frequencies for each meQTL for each pair of populations.} 
%     }

\generateFig{meqtlsallclocks}
{main_figs/Figure6.pdf}{1}
    {Genetic variants with different frequencies between ancestries associate with methylation levels at clock CpGs.}
    %Many methylation QTL (meQTL) for clock CpGs differ in frequency between populations.
    %Ancestry-specific genetic variants (meQTLs) drive Horvath clock error in African ancestry populations.
    %Ancestry-differentiated variants associate with methylation levels at clock CpGs.
    {{
        {\textbf{A}: Percentage of CpG sites in each methylation clock that exhibit differential methylation in whole blood between individuals of African vs. European ancestry.}
        {\textbf{B}: Methylation levels at some clock CpG sites significantly associate with error across MAGENTA individuals CpGs for each first generation clock are sorted based on the multiple-testing-corected significance of their association with clock error in MAGENTA.}
        {\textbf{C}: The allele frequency distribution of the single nucletide variants that disrupt CpG sites considered by the Horvath clock in 76,156 individuals from gnomAD (v3.0). Of the 353 clock CpG sites, 245 (69\%) have at least one variant, but nearly all the variants are very rare.
        }
        {\textbf{D}: Number of CpGs in each clock that are influenced by methylation quantitative trait loci (meQTLs). The Horvath clock contains 271 meQTL-associated CpGs, substantially more than the Hannum, EN, PhenoAge, or DunedinPACE clocks.
        }
        {\textbf{E}: Clock meQTL have significantly higher alelle frequency in individuals with African genetic ancestry from gnomAD than all other ancestry groups (median 0.068 for African vs. 0.004--0.046; p $<$ $3.85 \times 10^{-25}$).}
        {\textbf{F}: Frequency of meQTL variants in the MAGENTA cohorts. Non-differentiated meQTLs (left) have similar frequencies across cohorts, but African-differentiated meQTLs (right) are significantly more frequent in African Americans (AA) and Puerto Ricans (PR) compared to Peruvians (PER) and Cubans (CUB).}
        {\textbf{G}: Venn diagrams illustrating the overlap between CpGs that contribute to clock prediction error (``Error-associated'', blue) and those influenced by meQTLs (``meQTL affected'', pink) for the first-generation clocks.}
    }}

\section{Results}
Our primary analyses are based on genotyping and blood DNA methylation data from 621 individuals from the Americas with AD and non-demented controls collected by the MAGENTA study and $>$2,500 individuals from three replication cohorts. The MAGENTA study is focused on late-onset Alzheimer's and thus the average age of participants is 76 years old. Reflecting AD prevalence, the study is 68\% female. The individuals come from five cohorts, collected from the United States (White, African American, Cuban), Peru, and Puerto Rico (\cref{tab:cohorts}). We performed replication analyses on whole blood methylation data from African American individuals from the Grady Trauma Project (n = 422) and the GENOA study (n = 1,394), and White Swedish individuals (n = 729) that participated in the Northern Sweden Population Health Study (NSPHS). As described in detail in the Methods, to facilitate comparisons relevant to understand global differences in our study, we use a combination of geographic and race-based identifiers that are likely to best reflect underlying differences in genetic ancestry and admixture components. 

We apply a range of first-, second-, and third-generation methylation-based methylation clocks to these individuals. We then evaluate their accuracy in predicting chronological age, quantify whether they identify accelerated aging in individuals with AD, and explore genetic factors that may influence clock performance (\cref{fig:schematic}).

\schematic

\begin{table}[htbp]
    \centering
    \begin{tabular}{lccc}
        \toprule
        & \textbf{Alzheimer's (N=313)} & \textbf{Control (N=308)} & \textbf{Overall (N=621)} \\
        \midrule
        \textbf{COHORT} & & & \\
        African American & 98 (31.3\%) & 107 (34.7\%) & 205 (33.0\%) \\
        Cuban & 22 (7.0\%) & 21 (6.8\%) & 43 (6.9\%) \\
        White & 68 (21.7\%) & 65 (21.1\%) & 133 (21.4\%) \\
        Peruvian & 41 (13.1\%) & 41 (13.3\%) & 82 (13.2\%) \\
        Puerto Rican & 84 (26.8\%) & 74 (24.0\%) & 158 (25.4\%) \\
        \midrule
        \textbf{SEX} & & & \\
        Female & 206 (65.8\%) & 213 (69.2\%) & 419 (67.5\%) \\
        Male & 107 (34.2\%) & 95 (30.8\%) & 202 (32.5\%) \\
        \bottomrule
    \end{tabular}
    \caption{Demographics of the MAGENTA study cohorts.}
    \label{tab:cohorts}
\end{table}

\subsection{Methylation clock accuracy is lower in cohorts with substantial African ancestry}
To test whether current methylation clocks are able to predict age accurately in diverse, genetically admixed groups, we evaluated age predictions for the control individuals in the MAGENTA study. We first analyzed the widely used Horvath clock, which was trained on data from several tissues and cell types to accurately predict age across the lifespan using methylation levels at 353 CpG sites.

Age predicted from DNA methylation (DNAm age) in the White cohort using the Horvath clock was strongly correlated with chronological age (Pearson \textit{r} = 0.72) (\textbf{Figure 2A}). While this correlation is lower than reported in the original study ($>$0.9), it is consistent with previous studies of older individuals \citep{horvath_dna_2013, marioni_dna_2015}.

\correlationsancestry

In comparison to the White cohort, the correlation between DNAm age and chronological age was significantly lower for Puerto Ricans (\textit{r} = 0.45, \textit{p} = 0.007, \textit{n} = 74) and African Americans (\textit{r} = 0.51, \textit{p} = 0.016, \textit{n} = 107) (\textbf{Figure 2B}). The correlations for Cubans (\textit{r} = 0.68, \textit{ p} = 0.385, \textit{n} = 21) and Peruvians (\textit{r} = 0.72, \textit{p} = 0.52, \textit{n} = 41) were similar to the White cohort. These differences in clock accuracy were not driven by systematic differences in the age distributions for the cohorts (\textbf{Supplementary Figure 1}). The methylation clocks were also more accurate in the White cohort than the African American and Puerto Rican cohorts in terms of median absolute error (MAE) (\textbf{Supplementary Table 1}). Combining cases and controls did not qualitatively change the performance relationships; the accuracy remained lower for African Americans and Puerto Ricans relative to the Whites (\textbf{Supplementary Figure 3}).

The two cohorts with lower correlations come from regions where individuals often have substantial amounts of African ancestry. To explore if admixture levels associated with the accuracy of the Horvath clock in predicting age, we estimated the global proportions of African (YRI), European (CEU), and American (PEL) ancestries in each individual from the MAGENTA cohort using reference groups from the 1000 Genomes Project \citep{the_1000_genomes_project_consortium_global_2015}. 

Methylation clock accuracy was lowest for the cohorts with substantial African ancestry: African Americans (median 85\% African) and Puerto Ricans (median 15\% African). In contrast, the clocks performed similarly to the White cohort in groups with the lowest African ancestry: Cubans (6\% African) and Peruvians (2\% African) (\textbf{Figure 2C}). We regressed the Horvath clock error on the proportions of African ancestry in the genomes of the MAGENTA individuals, adjusting for chronological age. The proportion of African ancestry is significantly associated with increased Horvath clock error (p = 0.039), with an estimate of 1.46 years more error for 100\% African ancestry compared to no African ancestry.

\subsection{Lower methylation clock accuracy in African ancestry individuals replicates in multiple cohorts}
To evaluate if the significantly lower accuracy of the Horvath clock age predictions on individuals with substantial African ancestry holds outside of MAGENTA, we analyzed three independent whole blood methylation datasets. Two focused on African American individuals---the Grady Trauma Project (n = 422) \citep{katrinli_epigenome-wide_2024} and the GENOA study (n = 1,394) \citep{shang_meqtl_2023}---and one focused on White Swedish individuals (n = 729) \citep{johansson_continuous_2013}. Each study sampled a wide range of chronological ages, including many older individuals. To enable direct comparison with MAGENTA, we first focused on older individuals ($\geq 55$ years old). As observed in MAGENTA, the Horvath clock had lower accuracy for the African American cohorts (\textbf{Figure 3; bottom row}; Grady: \textit{r} = 0.57, GENOA: \textit{r} = 0.67) and higher accuracy for the White Swedish cohort (\textit{r} = 0.80). We then expanded the analysis to all individuals in each cohort. As expected, the overall accuracy of the age predictions improved, but the disparity in performance between cohorts remained (\textbf{Figure 3; top row}; Grady: \textit{r} = 0.88, GENOA: \textit{r} = 0.85, Swedish: \textit{r} = 0.97).

\magentareplication

%The MAE for the Grady Project individuals was even higher than that of the control MAGENTA African Americans (Grady Project ; MAGENTA AA MAE = 5.38), and the correlation between predicted and chronological age was virtually the same (Grady Project ; MAGENTA AA R = 0.58). The GENOA study individuals however had a lower MAE (4.12) relative to the Grady Project individuals and the MAGENTA African Americans and Puerto Ricans. Finally, the correlation between predicted and chronological age for these individuals also higher than the aforementioned cohorts (R = 0.67), but still lower than that of the MAGENTA White individuals (R = 0.72), despite the stark difference in sample size between these cohorts.

%Because the tissues sampled from these individuals are whole blood, and there is an older age range covered, clock accuracy for these individuals can be compared to the accuracy observed in the MAGENTA cohorts. To enforce this direct comparison, we restricted our analyses to those older individuals. Overall, 

%As a final sanity check, we gathered publicly available methylation data generated from whole blood samples of White Swedish individuals. We restricted our efforts to older individuals ($\geq 55$ years old), much in the same way that we performed the Grady Project and GENOA analyses, thereby making the results comparable to those obtained in the MAGENTA cohorts. We observed a much smaller MAE for the White Swedish individuals in comparison to all individuals with substantial African ancestry evaluated herein (3.73). The correlation between predicted and chronological age was also much higher, and comparable to that observed in the MAGENTA White individuals (R = 0.79) (\textit{Figure 3}. 

A full summary of the clock accuracy statistics, including MAE, maximum absolute error, correlation, and sample size for each cohort are included in \textbf{Supplementary Table 2}. Taken together, these results show that the reduced clock accuracy for individuals with substantial African ancestry is not limited to the cohorts in the MAGENTA study, but rather is a general pattern in the portability of methylation clocks.


\subsection{Accuracy of age prediction on admixed individuals varies across methylation clocks}
To investigate the performance of other methylation clocks at predicting chronological age in admixed individuals, we selected several additional publicly available open-source clocks. We considered two other ``first-generation'' clocks that were trained to predict chronological age: the Hannum clock \citep{hannum_genome-wide_2013} and a model developed by \citealp{zhang_improved_2019} that used large datasets for training and achieved substantially higher performance than previous age predictors. We hereafter refer to this elastic net model as ``Zhang\_EN''.

Both models achieved higher correlations with chronological age than the Horvath clock across the cohorts in the MAGENTA study. For example, Hannum has a correlation of 0.74 in the White cohort, and consistent with previous evaluations, Zhang\_EN has a correlation of 0.88. These relative performance trends held across cohorts, but again, the cohorts with substantial African ancestry,  African Americans and Puerto Ricans, consistently had the lowest age correlations for each clock (\textbf{Figure 4}). 

Next, we evaluated the PhenoAge clock, a ``second-generation'' clock that is trained on biomarkers of frailty and physiological deterioration \citep{levine_epigenetic_2018}. The correlation between DNAm age and chronological age was lower for this clock in comparison to the other methylation clocks (\textit{r} = 0.53 in the White cohort). This is likely due to the fact that this clock was not trained to predict age directly, but rather markers of aging. This clock did not show as substantial a difference in performance between cohorts as seen for the Horvath clock, but the African American and Puerto Rican individuals again had the lowest correlation of all cohorts. 

Overall, these results demonstrate that current methylation clocks vary in the correlation of their predicted DNAm age with chronological age in genetically admixed cohorts. The clocks are also consistently the least accurate in predicting age in cohorts with substantial proportions of African ancestry.

\accuracymultipleclocks

\subsection{Principal component versions of the methylation clocks also have lower age prediction accuracy for genetically admixed individuals}
Principal component (PC) versions of many methylation clocks have been developed with the aims of reducing noise and improving replicability and generalizability of age predictions \citep{higgins-chen_computational_2022}. Thus, we evalauted the accuracy and stability of PC methylation clocks across genetic ancestries. Overall, the PC version of the Horvath clock did not result in consistent improvement in age prediction accuracy or generalization across MAGENTA cohorts (\textbf{Supplementary Figures 4 and 5}). The lower accuracy for age prediction in individuals of substantial African ancestry were present for the PC versions of the clocks in the replication cohorts, just as in the MAGENTA cohorts (\textbf{Supplementary Figure 6}). Taken together, these results indicate that PC-based versions of the existing methylation clocks do not provide generalizable age predictors for individuals of diverse genetic ancestries.


\subsection{Most methylation clocks do not identify accelerated aging in admixed Alzheimer's cohorts}
DNAm age has been proposed as a promising biomarker and predictive tool for age-related disease risk, particularly because of associations between accelerated DNAm age (compared to chronological age) and the presence of diseases such as coronary heart disease, Parkinson's disease, and AD. However, these results have largely been observed in European-ancestry cohorts.

To evaluate the ability of methylation clocks to identify accelerated aging and risk for age-related disease in diverse, genetically admixed individuals, we quantified the association of methylation age acceleration with AD status in cohorts from the MAGENTA study. In addition to the clocks tested in the previous section, we also included a ``third generation'' clock, DunedinPACE, that aims to predict the pace of aging as measured by change in biomarkers over time from methylation data, rather than age itself \citep{belsky_dunedinpace_2022}. 

The cell type composition in blood is known to change with age, which if not accounted for, can confound age acceleration estimates \citep{jaiswal_clonal_2019}. Thus, we focused on intrinsic age acceleration estimates computed using established algorithms to correct for cell type composition. 

To establish a baseline for this analysis, we tested whether individuals with AD in the White cohort show significantly greater age acceleration than non-demented controls. As expected from previous studies \citep{levine_epigenetic_2015, levine_epigenetic_2018}, AD cases have modest but significantly greater age acceleration as measured by the Horvath clock than controls (\textbf{Supplementary Figure 7}; median 1.7 vs. 1.5 years, \textit{p} = 0.041). For each of the other clocks, AD cases had higher median age acceleration than controls (\textbf{Figure 5B}), though the differences only reached statistical significance for the DunedinPACE clock (median 1.09 vs. 1.07, \textit{p} = 0.044).

Having established that previously reported age acceleration in AD was detectable in the MAGENTA White cohort, we evaluated whether the clocks found accelerated aging in the admixed AD cohorts. Focusing first on the Horvath clock, we observed inconsistent relationships between age acceleration and AD status. None of the admixed cohorts showed a significant difference, and controls even had higher median age acceleration in Peruvians and Cubans  (\textbf{Figure 5A}). Across the other clocks, none consistently identified greater age acceleration in AD cases across all populations (\textbf{Figure 5B}). Among the first- and second-generation clocks, only PhenoAge demonstrated a significant ability to differentiate AD cases from controls in any of the non-European ancestry groups, specifically in African American individuals (p = 0.008), which were included in its training set. Cubans consistently showed greater age acceleration in controls rather than cases, while none of the other admixed cohorts even had consistent directions of effect across methods. While sample size in these cohorts is relatively small, our power calculations suggest that low statistical power is unlikely to explain the lack of signal (see Discussion).

DunedinPACE stood out in the evaluation, as it identified significantly greater aging in AD cases compared to controls in the White (p = 0.044), African American (p = 0.0019), and Puerto Rican (p = 0.0090) cohorts using its ``pace of aging'' metric. However, no significant differences were found for Cubans (p = 0.26) or Peruvians (p = 0.81).

\alzheimersanalysis


\subsection{Combining predictions across methylation clocks does not improve their performance}
Inspired by recent work on the ensembling of PRS to better predict disease risk from genetic variation across populations \citep{monti_evaluation_2024}, we evaluated whether combining age predictions could lead to greater accuracy in age prediction and AD risk prediction in the admixed cohorts. To investigate this, we averaged the age predictions for each individual in the MAGENTA study across five methylation clocks: Horvath, Hannum, Zhang\_EN, Zhang\_BLUP, and PhenoAge clocks. (The Zhang\_BLUP is a variation of the Zhang\_EN clock that does not use strong regularization.) 

The ensemble method's DNAm age prediction is more strongly correlated with chronological age in comparison to the Horvath and PhenoAge clocks, but it did not improve upon the best predictors (Zhang clocks) across populations (\textbf{Supplementary Figure 8A}).

We next evaluated whether the ensemble intrinsic age acceleration estimates would associate more strongly with AD disease status relative to the standalone methylation clocks. Following the same evaluation framework as for the individual clocks, we found only one significant difference in age acceleration. Cuban control individuals had significantly \textit{lower} age acceleration than AD cases (\textbf{Supplementary Figure 8B}). Thus, a simple ensemble does not lead to stronger performance at either task in admixed cohorts.

\subsection{Many clock CpGs are differentially methylated between European and African ancestry individuals}
Our results so far demonstrate that existing methylation clocks do not perform consistently across genetically admixed individuals. We now shift our attention to investigating possible mechanisms underlying this lack of generalization. We hypothesized the differential methylation levels between genetic ancestries at clock CpGs could contribute to the lower performance of methylation clocks in admixed individuals. To test for this pattern, we analyzed whole blood methylation data from the EWAS Hub for 306 individuals of African (AFR) and 300 of European (EUR) ancestry. We fitted a linear regression model to the whole blood methylation profiles for these two sets of individuals, including covariates for sex and chronological age. We called CpG sites as differentially methylated in AFR vs. EUR individuals controlling the false discovery rate at 5\% with the Benjamini-Hochberg correction. We identified 73,819 differentially methylated CpG sites and intersected them with CpGs included in each clock.

All of the clocks considered had a substantial proportion of their CpGs differentially methylated in AFR vs EUR individuals (\textbf{Figure 6A}). The PhenoAge and Horvath clocks had the highest fractions with 24.8\% (127/513) and 23.8\% (84/353) of their CpG sites differentially methylated, respectively. The other clocks also had substantial fractions of their CpGs differentially methylated: Zhang EN clock (12.6\%, 65/514), DunedinPACE (9.2\%, 16/173), and the Hannum clock (7\%, 5/71). This suggests  systematic differences in methylation patterns between genetic ancestries at CpG sites considered in methylation clocks are common.

%\clockproportions

\subsection{Methylation levels at many CpGs significantly associate with clock error}
The presence of differentially methylated CpGs in a clock does not necessarily indicate a barrier to accurate prediction across ancestries. For example, the differential methylation could reflect true differences in environmental exposures between the groups that influence aging. To identify CpGs that could drive the lower accuracy for individuals with substantial African ancestry in their genomes, we tested whether methylation levels at individual clock CpGs are associated with increased clock error across all MAGENTA individuals. The number of CpG sites with methylation levels significantly associated (after Benjamini-Hochberg multiple testing correction) with clock error in the MAGENTA individuals varied across the clocks (\textbf{Figure 6B}). The methylation levels at 56 of the 353 Horvath clock CpGs significantly associated with increased clock error, in contrast to 19 of the 71 Hannum clock CpGs. Finally, the methylation levels at 52 of the 514 Zhang EN clock CpGs signficantly associated with increased clock error in MAGENTA individuals. Because the PhenoAge and DunedinPACE clocks are second and third generation clocks, respectively, and therefore do not aim to predict age directly but rather a composite biological age, we did not include these two clocks in this analysis. 

\subsection{Do genetic factors contribute to lower methylation clock performance in admixed populations?}

Given that many CpGs sites used in clocks have different methylation levels between genetic ancestries, we sought to investigate potential mechanisms underlying the differenes in methylation and decreased performance of some clocks. Environmental differences between populations likely contribute to the decreased performance of the clocks; however, we do not have comprehensive data on differences in environmental exposures between the cohorts. Moreover, several observations suggest that genetic factors also play a substantial role in clock generalizability. First, genetic ancestry explains a substantial fraction of methylation differences between human populations~\citep{galanter_differential_2017}. Furthermore, the particularly large decrease in performance for individuals with a substantial fraction of African ancestry in different populations suggests that genetic divergence rather than the presence of environmental differences alone is responsible. 

Thus, we now focus on evaluating potential mechanisms by which genetic variation could reduce clock accuracy across human groups. First, we evaluate whether genetic variation in human populations disupts the potential for methylation at CpG sites. We then explore how often genetic variants that associate with methylation levels (meQTL) influence clock CpGs and differ in frequency between genetic ancestries.

\subsection{Many methylation clock CpGs are disrupted by genetic variants, but the variants are extremely low frequency}

We first quantified how often a genetic variant segregating in human populations disrupted a CpG site included in a clock. This scenario could lead to errors if the variant has different frequencies across cohorts given the loss of potential for methylation and the ability of the site to contribute to the age prediction. 

Of the 353 CpG sites considered in the Horvath clock, 245 (69\%) have at least one disruptive genetic variant observed in at least one individual in the gnomAD database of variants identified in a cohort of 76,156, including thousands of individuals of non-European ancestry. While this may seem like a substantial concern, these variants are extremely rare both globally and stratifying by ancestry (\textbf{Figure 6C; Supplementary Figure 9}). The average frequency is 0.0001, with the most common case being a variant observed in just one individual. Only one clock CpG disrupting variant had a frequency greater than 1\%. 

The genetic variant frequency patterns are similar for the other clocks (\textbf{Supplementary Figure 10}). Of the 71 CpG sites considered in the Hannum clock, 54 (76\%) have at least one disruptive variant from the gnomAD database, but only one of the variants has a frequency greater than 1\%. For the 514 CpG sites that make up the Zhang19\_EN clock, 403 (78\%) are disrupted by at least one gnomAD variant, and five of the variants in clock CpGs have a frequency greater than 1\%. The PhenoAge clock has 381 of its 513 (74\%) disrupted by a gnomAD variant, but only two of these variants have an allele frequency greater than 1\%. Finally, in the case of DunedinPACE clock and its 173 CpG sites, 158 (91\%) are disrupted by at least one gnomAD variant. However, only one of these variants has an allele frequency greater than 1\%. Thus, genetic variation in clock CpG sites themselves is unlikely to be the main cause of the lack of generalization of the methylation clocks.


\subsection{Common methylation QTL influence clock CpGs} 
We next assessed the prevalence of meQTLs, genetic variants that associate with clock CpG site methylation, another potential modifier of methylation levels that could lead to spurious DNAm age predictions across individuals. We gathered three sets of meQTLs from Europeans, South Asians, and African Americans (Methods). We intersected the CpGs associated with the meQTLs with clock CpG sites. 

Out of the 353 CpGs included in the Horvath clock, 271 (77\%) had at least one meQTL. Overall, a total of 29,033 unique variants associated with methylation levels at Horvath clock CpGs. In contrast to CpG disrupting variants, the meQTL had an average allele frequency of 0.19, and 26,500 were common, as determined from gnomAD version 4.1 (global allele frequencies $\geq$1\%). However, only a small proportion of CpG sites in the other first-generation clocks have meQTL (Hannum: 8\%, Zhang\_EN: 1\%). PhenoAge is similar with 7\% of its CpGs affected by at least one meQTL. Finally, DunedinPACE had no meQTLs affecting its 173 clock CpG sites (\textbf{Figure 6D}).

Thus, the clock with the largest decrease in performance in admixed cohorts (in terms of predicting chronological age and identifying age acceleration in AD) has the largest fraction of CpGs with meQTLs. In contrast, DunedinPACE, the best performing clock at identifying AD cases in the MAGENTA study, had no meQTLs. The three other clocks with intermediate performance in the admixed cohorts, all have meQTL for some CpGs, but much lower fraction than the Horvath clock.

\subsection{Clock CpG methylation QTLs vary in frequency across ancestries}
The presence of meQTL influencing clock CpGs does not necessarily indicate a problem for clock generalization. However, differences in the presence or frequency of meQTL that influence clock CpGs between genetic ancestries could lead to decreased DNAm age prediction accuracy (and therefore weaker associations with disease) in admixed cohorts. For example, if a clock is trained on a cohort without an meQTL, the learned weights for the CpG will not have accounted for the effects of the meQTL. To quantify whether differences in meQTL across genetic ancestries could influence methylation clocks, we analyzed the gnomAD allele frequencies for the 29,033 Horvath clock meQTL tag variants in multiple global populations. 

The Horvath clock meQTLs are at significantly higher frequencies in African ancestry populations (median 0.068) than in each of the eight other population groups considered (\textbf{Figure 6E}; 0.004--0.046, p $<$ $3.85 \times 10^{-25}$). There were also 2,328 meQTL that were only observed in African individuals.

To connect these results to individuals with recent admixture, like many in the MAGENTA study, we also tested whether the Horvath clock meQTL differed in frequency in local ancestry blocks of different origins in genetically admixed individuals from gnomAD. We used pre-computed local ancestry calls for 7,612 Latino/Admixed American individuals to compare allele frequencies for each meQTL in three ancestral backgrounds: African, Amerindigenous, and European. The clock CpG-affecting meQTLs were at higher frequencies in African local ancestry backgrounds (\textbf{Supplementary Figure 12}) relative to Amerindigenous and European backgrounds, consistent with our finding that these meQTLs are most frequent in African ancestry individuals at the global population level.

%\subsection{Susceptibility to meQTL varies across methylation clocks}
%We have thus far found that many of the methylation clocks are less accurate at predicting age for individuals of substantial AFR ancestry and that substantial proportions of the CpG sites in the clocks are differentially methylated in AFR individuals relative to EUR individuals. 

%We hypothesized that meQTL might affect the clock CpG sites, and that some of these meQTL could be genetic ancestry-specific. Differences in the presence or frequency of meQTL that influence clock CpGs between genetic ancestries could lead to decreased DNAm age prediction accuracy (and therefore weaker associations with disease) in diverse and admixed cohorts. For example, if a clock is trained on a cohort without an meQTL, the learned weights for the CpG will not have accounted for the effects of the meQTL. In addition, we hypothesized that population-specific environmental factors might play a role in affecting the clock CpG sites in a way that is not reflected in the training data for the clocks. Both of these factors could, at least partially, explain the lack of generalizability of existing methylation clocks to more diverse populations (\textbf{Figure 6B}).

%We first investigated the scenario of meQTL affecting the CpG sites in the methylation clocks. We gathered three sets of meQTLs from Europeans, South Asians, and African Americans and intersected the affected CpGs for the meQTLs with the Horvath clock CpG sites. We found a strong potential for meQTL to influence the Horvath clock (76\% of its CpGs have an meQTL) (\textbf{Figure 6C}), in stark contrast to the small proportion of affected clock CpG sites in the other first-generation clocks considered here (Hannum: 8\%, Zhang2019\_EN: 1\%). PhenoAge is similarly low, with 7\% of its CpGs affected by at least one meQTL. Finally, DunedinPACE had no meQTLs affecting its 173 clock CpG sites.

%Thus, the clock with the largest decrease in performance in admixed cohorts (in terms of predicting chronological age and identifying age acceleration in AD) has the most and largest fraction of meQTLs influencing its CpGs. On the opposite side of the spectrum, DunedinPACE, the best performing clock at identifying AD cases in the MAGENTA study, had no meQTLs. The three other clocks with intermediate performance in the admixed cohorts, all have meQTL for some CpGs, but much lower fraction than the Horvath clock.

\subsection{Differentiated meQTL are present in the MAGENTA cohort}
To explore if meQTL influence CpG methylation relevant to clock predictions in the MAGENTA individuals, we tested whether meQTLs described in the previous section are present in MAGENTA. For example, an meQTL (chr3:51639064\_A\_G) for a Horvath clock CpG with a reported beta of 1.34 and an allele frequency of 2\% in gnomAD (version 4.1) is present in MAGENTA individuals (study-wide frequency = 11\%; AA frequency = 20\%; PR frequency = 7.5\%; CUB and PER combined frequency = 6\%). In individuals with at least one copy of the variant, the methylation level of the corresponding Horvath clock CpG site was higher, as expected from the reported beta for the meQTL. The median methylation level of the clock CpG site for individuals with at least one copy of the variant was significantly higher than for individuals without the variant (median difference in methylation beta = 0.04; \textit{p} = $1.64 \times 10^{-7}$). Consistent with our hypothesis that this meQTL could contribute to Horvath clock error, individuals with this meQTL had 0.4 years higher median clock error. 

%To explore this pattern across clock CpGs and see if it could contribue to differences between admixed populations, we quantified the burden of clock CpG meQTL in each MAGENTA cohort. AA and PR individuals with at least one copy of a clock CpG meQTL, relative to the CUB and PER cohorts. We then interrogated whether this trend would remain if we analyzed clock CpG meQTL that are 5-fold more frequent in AFR ancestry individuals relative to the next most frequent ancestry in the gnomAD database. We consider this set of variants to be AFR-differentiated. As expected based on global ancestry proportions for the cohorts, we observed an even more reinforced version of the previous result.

Expanding on this example, we quantified whether differences in meQTL frequency across populations observed in other datasets are present in MAGENTA. The overall frequency of meQTL without substantial differentiation in large databases was similar between the MAGENTA admixed cohorts (Figure 6F, left panel; median variant frequencies: AA = 0.32, PR = 0.31, PER = 0.29, CUB = 0.30). However, as expected, the frequency of AFR-differentiated meQTL was higher in the MAGENTA African American (median = 0.21) and Puerto Rican (median = 0.07) cohorts compared to the Peruvian (median = 0.02) and Cuban (median = 0.04) cohorts (\textbf{Figure 6F, right panel}). Thus, the MAGENTA individuals carry many meQTL at different frequencies in African ancestry that influence clock CpG methylation.

\meqtlsallclocks

Finally, we tested whether the error-associated CpGs in the clocks were also affected by meQTL. Strikingly, of the 56 Horvath clock CpGs that were significantly associated with increased clock error in MAGENTA individuals, 42 were also affected by an meQTL, and nine were affected by African ancestry-differentiated meQTL. This pattern contrasts  with the Hannum and Zhang EN clocks, where only three and one of the error-associated CpGs were affected by an meQTL, respectively (\textbf{Figure 6G and Supplementary Figure 13}). 

%Taken together, these data show that the Horvath clock is the least accurate methylation clock when applied to genetically diverse individuals, with a substantial proportion of its CpGs affected by genetic variants, including those that are specific to African ancestry. These CpGs and their methylation levels are, in turn, associated with increased clock error in the MAGENTA individuals.

\clearpage